\documentclass[11pt]{article}
\usepackage{mathblatt}
\begin{document}
\blattkopf{Unabhängigkeit}{Lernblatt}
\verzeichniszeile{\verz{e1}{1 Unabhängigkeit prüfen: die Produktregel P(A ∩ B) = P(A) · P(B) an Anteilen oder absoluten Häufigkeiten prüfen (Tafel oder Anzahlen erst beschaffen (Nr.~9–11)} \verztrenn \verz{e2}{2 Rückwärts (Nr.~12–13)} \verztrenn \verz{e3}{3 Unabhängigkeit als Argument: die Ausgleichs-Fehlvorstellung widerlegen (Nr.~14–15)} \verztrenn \verz{abhaken}{Das kann ich}}
% Einheit 1 von 3 – Unabhängigkeit prüfen: die Produktregel P(A ∩ B) = P(A) · P(B) an Anteilen oder absoluten Häufigkeiten prüfen (Tafel oder Anzahlen erst beschaffen (bank/unabhaengigkeit/e1.jsonl, zusammenbau v0.1)
%% TODO Zeitmarke Einheit 1: Marken-Zeile nicht lesbar
%% TODO Zweigzeile Teil 1: was hier gelernt wird (Satz in Schülersprache aus dem Einheitstitel) – Einheit 1
\einheitenkopf[e1]{Einheit 1 von 3 · Unabhängigkeit prüfen: die Produktregel P(A ∩ B) = P(A) · P(B) an Anteilen oder absoluten Häufigkeiten prüfen (Tafel oder Anzahlen erst beschaffen}
\zweigzeile{Abitur GK}
\setcounter{aufgabe}{8}

%% TODO Titel: Ich-kann-Satz fehlt in der Bank (Platzhalter: Kettenname) – Nr. 9
%% TODO Anweisung: ein Satz über der ersten Teilaufgabe fehlt in der Bank – Nr. 9
\begin{aufgabe}{„Prüfen oder nutzen?“}
\begin{teile}
\teil Aufgabe: „Untersuche, ob die Ereignisse A und B stochastisch unabhängig sind.“ Was ist zu tun? \\ \kreuz{prüfen – die Produktregel anwenden} \\ \kreuz{nutzen – Unabhängigkeit ist vorausgesetzt, eine Gleichung aufstellen}
\end{teile}
\end{aufgabe}

%% TODO Titel: Ich-kann-Satz fehlt in der Bank (Platzhalter: Kettenname) – Nr. 10
%% TODO Anweisung: ein Satz über der ersten Teilaufgabe fehlt in der Bank – Nr. 10
\begin{aufgabe}{Unabhängigkeit prüfen}
%% TODO \rechenplatz{n}: Schrittzahl des Grundfalls fehlt in der Bank (nur bei fester Schreibform, 2.3 b)
\begin{teile}
\teil Eine Vierfeldertafel mit allen Feldern und Rändern liegt vor. Wo stehen die drei Zahlen der Produktregel? \\ \kreuz{alle drei stehen in der Tafel: zwei Ränder und ein Feld} \\ \kreuz{nur die Ränder stehen da, der Schnitt muss berechnet werden} \\ \kreuz{keine steht da, alle drei sind zu beschaffen}
\teil Die Tafel zeigt 100 Schüler nach den Ereignissen M (Mädchen) und F (kommt mit dem Fahrrad). Prüfe mit der Produktregel, ob M und F stochastisch unabhängig sind.

\vierfeldertafel{M}{F}{20,30,50,20,30,50,40,60,100}
\teil In einem Sportverein sind 60 Mitglieder, 36 davon Erwachsene. 20 Mitglieder spielen Volleyball, darunter 8 Jugendliche. Bestimme die fehlenden Anzahlen und prüfe, ob E (Erwachsener) und V (spielt Volleyball) stochastisch unabhängig sind.
\teil In einer Region traten 12\,400 Personen zur Fahrprüfung an, davon waren 3\,100 mindestens 30 Jahre alt. 9\,920 Prüflinge bestanden, darunter 7\,750 jüngere. Prüfe, ob der Anteil der Bestandenen unter den Älteren mit dem Gesamtanteil übereinstimmt, ob Ä (mindestens 30 Jahre) und S (bestanden) stochastisch unabhängig sind, und deute das Ergebnis \hfill (Abitur 2019 GK).
\teil Ein Paketdienst befördert 600 Pakete, 150 davon sind schwer. 90 der schweren Pakete gehen ins Ausland, insgesamt gehen 240 Pakete ins Ausland. Untersuche, ob der Anteil der Auslandspakete unter den schweren Paketen ebenso groß ist wie unter den nicht schweren \hfill (Abitur 2022 GK).
\teil Ein Glücksrad mit vier gleich großen Feldern 1, 2, 3, 4 wird zweimal gedreht. Ereignis A: die Summe der beiden Zahlen ist 5. Ereignis B: beim ersten Drehen erscheint die 4. Prüfe ohne Rechner, ob A und B stochastisch unabhängig sind \hfill (Abitur 2018 LK).
\teil Von den Kunden eines Versandhauses besitzen 70\,\% eine Kundenkarte (Ereignis K). Von den Karteninhabern bestellen 40\,\% online; 15\,\% aller Kunden sind Kunden ohne Karte, die online bestellen. O bezeichnet das Ereignis „bestellt online“. Begründe mithilfe der Ungleichung $0{,}4 \ne 0{,}7 \cdot 0{,}4 + 0{,}15$, dass K und O stochastisch abhängig sind \hfill (Abitur 2022 LK).
\end{teile}
\end{aufgabe}

%% TODO Titel: Ich-kann-Satz fehlt in der Bank (Platzhalter: Kettenname) – Nr. 11
%% TODO Anweisung: ein Satz über der ersten Teilaufgabe fehlt in der Bank – Nr. 11
\begin{aufgabe}{Unabhängigkeit prüfen – Fehler finden, Begründen, Anwendung, Darstellungswechsel}
\begin{teile}
\teil Die Tafel zeigt 100 Schüler nach M (Mädchen) und F (kommt mit dem Fahrrad). Lars sagt: „Es gibt 20 Mädchen mit Fahrrad, der Schnitt ist also nicht leer – M und F sind abhängig.“ Finde den Fehler.

\vierfeldertafel{M}{F}{20,30,50,20,30,50,40,60,100}
\teil Begründe, warum die Prüfung $P(A \cap B) = P(A) \cdot P(B)$ und die Prüfung $P_A(B) = P(B)$ zum selben Ergebnis führen.
\teil In einer Stadt mit 40\,000 Haushalten beziehen 12\,000 Ökostrom. 8\,000 Haushalte besitzen ein Elektroauto, darunter 3\,600 Ökostromkunden. Der Energieversorger will Elektroautobesitzer gezielt für Ökostrom werben. Prüfe, ob „Ökostrom“ und „Elektroauto“ stochastisch unabhängig sind, und sage, ob die gezielte Werbung Sinn hat.
\teil Die Tafel zeigt 100 Schüler nach M (Mädchen) und F (kommt mit dem Fahrrad). Übertrage sie in ein Baumdiagramm mit M in der ersten Stufe und lies an den Ästen der zweiten Stufe ab, ob M und F unabhängig sind.

\vierfeldertafel{M}{F}{20,30,50,20,30,50,40,60,100}
\end{teile}
\end{aufgabe}

\clearpage
% Einheit 2 von 3 – Rückwärts (bank/unabhaengigkeit/e2.jsonl, zusammenbau v0.1)
%% TODO Zeitmarke Einheit 2: Marken-Zeile nicht lesbar
%% TODO Zweigzeile Teil 1: was hier gelernt wird (Satz in Schülersprache aus dem Einheitstitel) – Einheit 2
\einheitenkopf[e2]{Einheit 2 von 3 · Rückwärts}
\zweigzeile{Abitur LK}
\setcounter{aufgabe}{11}

%% TODO Titel: Ich-kann-Satz fehlt in der Bank (Platzhalter: Kettenname) – Nr. 12
%% TODO Anweisung: ein Satz über der ersten Teilaufgabe fehlt in der Bank – Nr. 12
\begin{aufgabe}{Rückwärts}
%% TODO \rechenplatz{n}: Schrittzahl des Grundfalls fehlt in der Bank (nur bei fester Schreibform, 2.3 b)
\begin{teile}
\teil „Die Ereignisse A und B sind stochastisch unabhängig; $P(A) = 0{,}3$ und $P(A \cap B) = 0{,}12$. Berechne $P(B)$.“ Was ist zu tun? \\ \kreuz{prüfen – Produktregel anwenden} \\ \kreuz{nutzen – aus der Unabhängigkeit eine Gleichung gewinnen}
\teil Von den Besuchern eines Museums sind 60\,\% Erwachsene; insgesamt kaufen 35\,\% der Besucher einen Katalog. Welcher Anteil der Erwachsenen müsste einen Katalog kaufen, damit „Erwachsener“ und „Katalogkauf“ stochastisch unabhängig sind? Begründe ohne Rechnung.
\teil Die Vierfeldertafel zeigt Wahrscheinlichkeiten zu den Ereignissen A und B mit einem Parameter p. Für welchen Wert von p sind A und B stochastisch unabhängig? \hfill (Abitur 2021 LK)

\vierfeldertafel{A}{B}{p,3p,4p,0{,}2,0{,}8 - 4p,1 - 4p,p + 0{,}2,0{,}8 - p,1}
\teil Ein Glücksrad hat drei Felder mit den Zahlen 1, 2 und 3; die 1 und die 2 erscheinen je mit der Wahrscheinlichkeit p, die 3 mit $1 - 2p$. Es wird zweimal gedreht. Ereignis E: die erste Drehung zeigt die 3. Ereignis G: beide Drehungen zeigen dieselbe Zahl. E und G sind stochastisch unabhängig. Stelle eine Gleichung auf, aus der p berechnet werden kann, und löse sie \hfill (Abitur 2023 LK).
\teil Die Ereignisse A und B sind stochastisch unabhängig mit $P(A) = 0{,}4$ und $P(A \cap B) = 0{,}1$. Berechne $P(B)$. P(B) = \leerfeld
\teil Ein Gerät ist defekt, wenn mindestens einer der Fehler A oder B auftritt; die beiden Fehler treten stochastisch unabhängig voneinander auf. Ein Gerät ist mit der Wahrscheinlichkeit $0{,}3$ defekt, Fehler A tritt mit der Wahrscheinlichkeit $0{,}2$ auf. Berechne die Wahrscheinlichkeit für Fehler B \hfill (Abitur 2024 LK). P(B) = \leerfeld
\teil Für zwei stochastisch unabhängige Ereignisse A und B gilt $P(B) = P(A) + 0{,}3$ und $P(A \cap \overline{B}) = 0{,}1225$. Bestimme $P(A)$ \hfill (Abitur 2025 LK).
\end{teile}
\end{aufgabe}

%% TODO Titel: Ich-kann-Satz fehlt in der Bank (Platzhalter: Kettenname) – Nr. 13
%% TODO Anweisung: ein Satz über der ersten Teilaufgabe fehlt in der Bank – Nr. 13
\begin{aufgabe}{Rückwärts – Fehler finden, Begründen, Anwendung}
\begin{teile}
\teil A und B sind stochastisch unabhängig mit $P(A) = 0{,}5$ und $P(B) = 0{,}3$. Für $P(A \cap \overline{B})$ rechnet Mia: \rechnung{P(A \cap \overline{B}) &= 0{,}5 - 0{,}3 = 0{,}2} Finde den Fehler.
\teil A und B sind stochastisch unabhängig. Begründe, warum dann auch $P(A \cap \overline{B}) = P(A) \cdot P(\overline{B})$ gilt.
\teil Ein Rauchmelder soll bei Rauch mit mindestens $99\,\%$ Wahrscheinlichkeit anschlagen. Er hat zwei unabhängige Sensoren; der erste erkennt Rauch mit $0{,}9$. Wie zuverlässig muss der zweite Sensor mindestens sein?
\end{teile}
\end{aufgabe}

\clearpage
% Einheit 3 von 3 – Unabhängigkeit als Argument: die Ausgleichs-Fehlvorstellung widerlegen (bank/unabhaengigkeit/e3.jsonl, zusammenbau v0.1)
%% TODO Zeitmarke Einheit 3: Marken-Zeile nicht lesbar
%% TODO Zweigzeile Teil 1: was hier gelernt wird (Satz in Schülersprache aus dem Einheitstitel) – Einheit 3
\einheitenkopf[e3]{Einheit 3 von 3 · Unabhängigkeit als Argument: die Ausgleichs-Fehlvorstellung widerlegen}
\zweigzeile{Abitur GK}
\setcounter{aufgabe}{13}

%% TODO Titel: Ich-kann-Satz fehlt in der Bank (Platzhalter: Kettenname) – Nr. 14
%% TODO Anweisung: ein Satz über der ersten Teilaufgabe fehlt in der Bank – Nr. 14
\begin{aufgabe}{Als Argument}
%% TODO \rechenplatz{n}: Schrittzahl des Grundfalls fehlt in der Bank (nur bei fester Schreibform, 2.3 b)
\begin{teile}
\teil Eine Münze zeigt dreimal hintereinander Wappen. Wie wahrscheinlich ist beim nächsten Wurf Zahl? \\ \kreuz{$0{,}5$} \\ \kreuz{größer als $0{,}5$} \\ \kreuz{kleiner als $0{,}5$}
\teil Ein Glücksrad zeigt bei jeder Drehung mit $25\,\%$ Wahrscheinlichkeit Blau. Bei 80 Drehungen erschien Blau nur 12-mal. Mia folgert, dass bei den nächsten 80 Drehungen der Anteil von Blau deutlich über $25\,\%$ liegen muss. Beurteile die Aussage \hfill (Abitur 2017 LK).
\teil Bei einem Online-Quiz beantwortet Jan jeden Samstag 8 Fragen; jede Frage beantwortet er unabhängig von den anderen mit der Wahrscheinlichkeit $0{,}6$ richtig. Die Anzahl der richtigen Antworten an einem Samstag ist binomialverteilt. An drei Samstagen hintereinander hatte er jeweils 7 richtige. Jan meint, die Chance auf 7 richtige sei an diesem Samstag deshalb deutlich kleiner. Beurteile die Aussage \hfill (Abitur 2021 GK).
\end{teile}
\end{aufgabe}

%% TODO Titel: Ich-kann-Satz fehlt in der Bank (Platzhalter: Kettenname) – Nr. 15
%% TODO Anweisung: ein Satz über der ersten Teilaufgabe fehlt in der Bank – Nr. 15
\begin{aufgabe}{Als Argument – Fehler finden, Begründen}
\begin{teile}
\teil Nach 50 Münzwürfen mit nur 18-mal Wappen sagt Paul: „Das Gesetz der großen Zahlen sagt, dass sich die Häufigkeit bei der Hälfte einpendelt – also muss jetzt öfter Wappen kommen.“ Finde den Fehler in der Begründung.
\teil Begründe, warum die Wahrscheinlichkeit für Wappen beim nächsten Münzwurf nicht davon abhängt, was in den bisherigen Würfen gefallen ist.
\end{teile}
\end{aufgabe}

% Abhakseite „Das kann ich“ – Zone nur im Gesamt (\mitzone)
%% TODO Ich-kann-Sätze: Platzhalter wie in den Aufgabendateien
\begin{abhakseite}
\ifdefined\mitzone
\abhakgruppe{Kennst du schon}
\abhak{1}{Vierfeldertafel füllen (Ränder, Felder, bedingte Angaben)}
\abhak{2}{Bedingte Wahrscheinlichkeit als Quotient bilden}
\abhak{3}{Anteile aus absoluten Häufigkeiten bilden und kürzen}
\abhak{4}{Pfadterme mit Parameter im Baumdiagramm (Ast als Faktor, Pfad als Produkt, Randwahrscheinlichkeit als Pfadsumme)}
\abhak{5}{Lineare und quadratische Gleichungen aufstellen und lösen, Lösungen sieben (Nulllösung verwerfen)}
\abhak{6}{Pfadregeln und Gegenereignis mit und ohne Zurücklegen}
\abhak{7}{Vierfeldertafel füllen (Ränder, Felder, bedingte Angaben) – Fehler finden}
\abhak{8}{Vierfeldertafel füllen (Ränder, Felder, bedingte Angaben) – selbst rechnen}
\fi
\abhakgruppe{Einheit 1 · Unabhängigkeit prüfen: die Produktregel P(A ∩ B) = P(A) · P(B) an Anteilen oder absoluten Häufigkeiten prüfen (Tafel oder Anzahlen erst beschaffen}
\abhak{9}{„Prüfen oder nutzen?“}
\abhak{10}{Unabhängigkeit prüfen}
\abhak{11}{Unabhängigkeit prüfen – Fehler finden, Begründen, Anwendung, Darstellungswechsel}
\abhakgruppe{Einheit 2 · Rückwärts}
\abhak{12}{Rückwärts}
\abhak{13}{Rückwärts – Fehler finden, Begründen, Anwendung}
\abhakgruppe{Einheit 3 · Unabhängigkeit als Argument: die Ausgleichs-Fehlvorstellung widerlegen}
\abhak{14}{Als Argument}
\abhak{15}{Als Argument – Fehler finden, Begründen}
\end{abhakseite}

\end{document}
